\def\thoigian{2}
\section
[Bài tập cuối chương I]
{BÀI TẬP CUỐI CHƯƠNG I}

\subsection{Mục tiêu}

\paragraph{Về kiến thức, kĩ năng}

\begin{itemize}
	\item Hệ thống các kiến thức đã học trong chương và cung cấp một số bài tập có nội dung tổng hợp, liên kết các kiến thức, kĩ năng đã học trong chương.
\end{itemize}

\paragraph{Về năng lực}

\begin{itemize}
	\item Rèn luyện năng lực toán học, nói riêng là năng lực mô hình hoá toán học, năng lực tư duy và lập luận toán học.
	\item Bồi dưỡng hứng thú học tập, ý thức làm việc nhóm, ý thức tìm tòi, khám phá và sáng tạo cho HS. 
\end{itemize}

\paragraph{Về phẩm chất}

\begin{itemize}
	\item Góp phần giúp HS rèn luyện và phát triển các phẩm chất tốt đẹp (yêu nước, nhân ái, chăm chỉ, trung thực, trách nhiệm): 
	\begin{itemize}
		\item Tích cực phát biểu, xây dựng bài và tham gia các hoạt động nhóm; 
		\item Có ý thức tích cực tìm tòi, sáng tạo trong học tập; phát huy điểm mạnh, khắc phục các điểm yếu của bản thân. 
	\end{itemize}
\end{itemize}

\subsection{Thiết bị dạy học và học liệu}

\paragraph{Giáo viên}

\begin{itemize}
	\item Giáo án, bảng phụ, máy chiếu (nếu có), giấy A3, … 
\end{itemize}

\paragraph{Học sinh}

\begin{itemize}
	\item SGK, vở ghi, dụng cụ học tập, máy tính cầm tay.
\end{itemize}

\subsection{Tiến trình dạy học}

\danhsachtiet

\subsubsection
[Ôn tập lí thuyết. Các bài tập trắc nghiệm và một số bài tập tự luận]
{Ôn tập lí thuyết.\\Các bài tập trắc nghiệm và một số bài tập tự luận}

\paragraph{Hoạt động khởi động}

\muctieu

\begin{itemize}
	\item Nhắc lại toàn bộ lí thuyết của chương I.
\end{itemize}

\noidungsanpham

\begin{ontap}
	Vẽ sơ đồ tư duy tổng hợp lại toàn bộ lí thuyết chương I: 
	\begin{listEX}
		\item hệ phương trình bậc nhất hai ẩn và các phương pháp giải;
		\item giải bài toán bằng cách lập hệ phương trình bậc nhất hai ẩn.
	\end{listEX}
	%%%
	\doicot
	%%%
	Bài làm của HS.
\end{ontap}

\tochucthuchien

\begin{tcth}
	\buocmot
	\item GV chia lớp thành các nhóm theo tổ, HS hoạt động theo nhóm.
	\buochai
	\item HS thực hiện vẽ sơ đồ tư duy.
	\buocba
	\item Các nhóm sẽ trình bày sản phẩm của nhóm mình trên	bảng, các nhóm khác theo dõi,
	nhận xét.
	\item Các HS khác theo dõi bài làm, nhận xét và góp ý.
	\buocbon
	\item GV tổng kết.
\end{tcth}

\paragraph{Hoạt động luyện tập}

\muctieu

\begin{itemize}
	\item Củng cố giải hệ phương trình bậc nhất hai ẩn bằng phương pháp thế, phương pháp cộng đại số hoặc sử dụng máy tính cầm tay.
\end{itemize}

\noidungsanpham

\begin{baitap}[1.19]
	Cặp số nào sau đây là nghiệm của hệ phương trình $\heva{&5x+7y=-1\\&3x+2y=-5}$?
	\choice
	{$(-1;1)$}
	{$(-3;2)$}
	{$(2;-3)$}
	{$(5;-5)$}
	%%%
	\doicot
	%%%
	Đáp án là B.
\end{baitap}

\begin{baitap}[1.20]
	Trên mặt phẳng toạ độ $Oxy$, cho các điểm $A(1; 2)$, $B(5; 6)$, $C(2; 3)$, $D(-1; -1)$.\\
	Đường thẳng $4x - 3y = -1$ đi qua hai điểm nào trong các điểm đã cho?
	\choice
	{$A$ và $B$}
	{$B$ và $C$}
	{$C$ và $D$}
	{$D$ và $A$}
	%%%
	\doicot
	%%%
	Đáp án là C.
\end{baitap}

\begin{baitap}[1.21]
	Hệ phương trình $\heva{&1{,}5x - 0{,}6y = 0,3 \\ &-2x + y = -2}$.
	\choice
	{có nghiệm là $(0; -0,5)$}
	{có nghiệm là $(1; 0)$}
	{có nghiệm là $(-3; -8)$}
	{vô nghiệm}
	%%%
	\doicot
	%%%
	Đáp án là C.
\end{baitap}

\begin{baitap}[1.22]
	Hệ phương trình $\heva{&0,6x + 0,3y = 1,8 \\ &2x + y = -6}$.
	\choice
	{có một nghiệm}
	{vô nghiệm}
	{có vô số nghiệm}
	{có hai nghiệm}
	%%%
	\doicot
	%%%
	Đáp án là B.
\end{baitap}

\begin{baitap}[1.23]
	Giải các hệ phương trình:
	\begin{listEX}
		\item $\heva{&2x + 5y = 10 \\ &\dfrac{2}{5}x + y = 1;}$
		\item $\heva{&0,2x + 0,1y = 0,3 \\ &3x + y = 5;}$
		\item $\heva{&\dfrac{3}{2}x - y = \dfrac{1}{2} \\ &6x - 4y = 2.}$
	\end{listEX}
	%%%
	\doicot
	%%%
	\huongdan
	\begin{listEX}
		\item Vô nghiệm;
		\item $(2;-1)$;
		\item Vô số nghiệm: $\left(x;\dfrac{3}{2}x-\dfrac{1}{2}\right)$, $x\in\mathbb{R}$.
	\end{listEX}
\end{baitap}

\begin{baitap}[1.24]
	Giải các hệ phương trình:
	\begin{listEX}
		\item $\heva{&0,5x + 2y = -2,5 \\ &0,7x - 3y = 8,1;}$
		\item $\heva{&5x - 3y = -2 \\ &14x + 8y = 19;}$
		\item $\heva{&2(x - 2) + 3(1 + y) = -2 \\ &3(x - 2) - 2(1 + y) = -3.}$
	\end{listEX}
	%%%
	\doicot
	%%%
	\huongdan
	\begin{listEX}
		\item $(3;-2)$;
		\item $(0{,}5;1{,}5)$;
		\item $(1;-1)$.
	\end{listEX}
\end{baitap}

\tochucthuchien

\begin{tcth}[baitap={từ 1.19 đến 1.22}]
	\buocmot
	\item GV tổ chức cho HS làm cá nhân.
	\buochai
	\item HS chuẩn bị câu trả lời.
	\buocba
	\item HS giơ tay phát biểu.
	\item Các HS khác nhận xét bài làm của bạn.
	\buocbon
	\item GV nhận xét, góp ý nếu cần.
\end{tcth}

\begin{tcth}[baitap={1.23 và 1.24}]
	\buocmot
	\item GV tổ chức cho HS làm cá nhân.
	\buochai
	\item HS làm bài vào vở.
	\item GV mời HS lên bảng làm bài.
	\buocba
	\item Các HS khác nhận xét bài làm của bạn.
	\buocbon
	\item GV nhận xét, góp ý nếu cần.
\end{tcth}

\paragraph{Tổng kết và hướng dẫn công việc ở nhà}

GV tổng kết lại nội dung bài học và dặn dò công việc ở nhà cho HS:

\begin{itemize}
	\item GV tổng kết lại các kiến thức trọng tâm của bài học: Các cách giải hệ hai phương trình bậc nhất hai ẩn.
	\item Nhắc HS về nhà ôn tập các nội dung đã học.
\end{itemize}

\subsubsection{Các bài tập tự luận cuối chương}

\paragraph{Hoạt động luyện tập}

\muctieu

\begin{itemize}
	\item Củng cố kĩ năng giải bài toán bằng cách lập hệ phương trình.
\end{itemize}

\noidungsanpham

\begin{baitap}[1.25]
	Tìm số tự nhiên $n$ có hai chữ số, biết rằng nếu viết thêm chữ số 3 vào giữa hai chữ số của số $n$ thì được một số lớn hơn số $2n$ là $585$ đơn vị, và nếu viết hai chữ số của số $n$ theo thứ tự ngược lại thì được một số nhỏ hơn số $n$ là $18$ đơn vị.
	%%%
	\doicot
	%%%
	\huongdan\\
	$n = 75$
\end{baitap}

\begin{baitap}[1.26]
	Trên cánh đồng có diện tích $160$ ha của một đơn vị sản xuất, người ta dành $60$ ha để cấy thí điểm giống lúa mới, còn lại vẫn cấy giống lúa cũ. Khi thu hoạch, đầu tiên người ta gặt $8$ ha giống lúa cũ và $7$ ha giống lúa mới để đối chứng. Kết quả là $7$ ha giống lúa mới cho thu hoạch nhiều hơn $8$ ha giống lúa cũ là $2$ tấn thóc. Biết rằng tổng số thóc (cả hai giống) thu hoạch cả vụ trên $160$ ha là $860$ tấn. Hỏi năng suất của mỗi giống lúa trên $1$ ha là bao nhiêu tấn thóc?
	%%%
	\doicot
	%%%
	\huongdan\\
	Năng suất của giống lúa cũ là 5 tấn/ha; của giống lúa mới là 6 tấn/ha.
\end{baitap}

\begin{baitap}[1.27]
	Hai vật chuyển động đều trên một đường tròn đường kính $20$ cm, xuất phát cùng một lúc, từ cùng một điểm. Nếu chuyển động ngược chiều thì cứ sau $4$ giây chúng lại gặp nhau. Nếu chuyển động cùng chiều thì cứ sau $20$ giây chúng lại gặp nhau. Tính vận tốc (cm/s) của mỗi vật.
	%%%
	\doicot
	%%%
	\huongdan\\
	$3\pi$ cm/s và $2\pi$ cm/s.
\end{baitap}

\begin{baitap}[1.28]
	Một người mua hai loại hàng và phải trả tổng cộng là $21{,}7$ triệu đồng, kể cả thuế giá trị gia tăng (VAT) với mức $10 \%$ đối với loại hàng thứ nhất và $8 \%$ đối với loại hàng thứ hai. Nếu thuế VAT là $9 \%$ đối với cả hai loại hàng thì người đó phải trả tổng cộng $21{,}8$ triệu đồng. Hỏi nếu không kể thuế VAT thì người đó phải trả bao nhiêu tiền cho mỗi loại hàng?
	%%%
	\doicot
	%%%
	\huongdan\\
	Loại hàng thứ nhất: 5 triệu đồng; Loại hàng thứ hai: 15 triệu đồng.
\end{baitap}

\tochucthuchien

\begin{tcth}[baitap={từ 1.25 đến 1.28}]
	\buocmot
	\item GV tổ chức cho HS làm cá nhân.
	\buochai
	\item HS làm bài vào vở.
	\item GV mời HS lên bảng làm bài.
	\buocba
	\item Các HS khác nhận xét bài làm của bạn.
	\buocbon
	\item GV nhận xét, góp ý nếu cần.
\end{tcth}

\paragraph{Tổng kết và hướng dẫn công việc ở nhà}

GV tổng kết lại nội dung bài học và dặn dò công việc ở nhà cho HS:
\begin{itemize}
	\item GV tổng kết lại các kiến thức trọng tâm của bài học: cách giải bài toán bằng cách lập hệ phương trình.
	\item Nhắc HS về nhà ôn tập các nội dung đã học.
\end{itemize}